% Research manuscript.  No Lean execution or original-point assertion.
\subsection{An actual ordinary cusp form with auxiliary order-$269$ inertia}
\label{sec:cubic-auxiliary-order269-seed}

Put $p=269$, $K=\mathbf Q(\sqrt{-3})$, $\eta=\chi_{-3}$,
and $\psi=\chi_{-4}$.  Dirichlet characters in formulas for modular
forms are evaluated on rational integers.  Their Galois avatars
$\operatorname{DirG}(\xi)$ are evaluated using arithmetic Frobenius.
Write $c_p$ for the cyclotomic character with
$c_p(\operatorname{Frob}^{\rm ar}_r)=r$.
The attached representation $R_g$ uses the native geometric convention
\[
 \det(1-XR_g(\operatorname{Frob}^{\rm geo}_r))
 =1-a_r(g)X+\xi(r)r^{k-1}X^2
 \quad(g\in S_k(N,\xi),\ r\nmid Np).
 \tag{G119.1}
\]
Thus $\det R_g=\operatorname{DirG}(\xi)^{-1}c_p^{1-k}$.
This inverse is material for the character of order $p$ below.

\begin{theorem}[AC-G119-auxiliary-seed: an actual auxiliary ordinary input]
\label{prop:cubic-auxiliary-order269-seed}
Let $q=3229$, and let $\theta$ be the primitive Dirichlet character
modulo $q$ specified by $\theta(6)=\zeta_{269}$.
For every odd integer $k>2$ there is a normalized cuspidal newform
\[
 g\in S_k(3q,\eta\theta)
 \tag{G119.2}
\]
of exact conductor $3q$, and a place above $269$ of its coefficient
field containing the values of $\theta$, such that
\[
 a_r(g)\equiv 1+\eta(r)r^{k-1}\pmod{\varpi}
 \quad(r\nmid3q),\qquad
 a_3(g)\equiv a_q(g)\equiv1\pmod{\varpi}.
 \tag{G119.3}
\]
In particular $a_p(g)\equiv1$, so $g$ is $p$-ordinary.

Take concretely $e=267$ and $k=2e+1=535$; the same conclusions below
hold for any positive odd $e$ with $268\nmid2e$.
Let $\theta^{1/2}=\theta^{135}$ and let $f_0$ be the primitive
automorphic twist
\[
 f_0=g\otimes(\eta\psi\theta^{-1/2}).
 \tag{G119.4}
\]
Then $f_0$ has weight $2e+1$, nebentype $\eta$, exact conductor
$48q^2$, and is ordinary at $269$.
Its normalized, absolutely irreducible representation
\[
 V=R_{f_0}c_p^e
   =R_g\,\eta\psi\,\Theta^{1/2}c_p^e,
 \qquad \Theta=\operatorname{DirG}(\theta),
 \tag{G119.5}
\]
satisfies
\[
 \det V=\eta,\qquad
 \overline V^{\rm ss}
   =\eta\overline\psi\,\overline c_p^{\,e}
       \oplus\overline\psi^{-1}\overline c_p^{-e}.
 \tag{G119.6}
\]
At $p$ it is crystalline ordinary with inertia characters
$\psi c_p^e$ and $\psi c_p^{-e}$.
At $q$ its exact finite inertia is
$\Theta^{1/2}\oplus\Theta^{-1/2}$, with $N=0$.
At $2$ it has scalar inertia $\psi$; at $3$ its inertia is
$1\oplus\eta$; the monodromy is zero at both places, and all other
places different from $p$ are unramified.

Let $T/\mathbf Q$ be the field cut out by $\theta$.
It is totally real cyclic of degree $269$, ramified only at $q$.
Restriction to $G_T$ kills the new $q$-inertia and preserves the
absolute irreducibility of $V$.
For this explicit $q$, $\theta(p)=\zeta_{269}^{175}$, so the prime
$p$ is unramified of residue degree $269$ in $T$, rather than split.

The assertion supplies an actual ordinary auxiliary-level rank-two
input.  It does not assert that $V$ is an original deformation point:
its nontrivial $p$-power inertia at $q$ violates the original
residual-kernel Kill condition there.  A unitary host, its integral
Hecke modules, and its numerical and propagation hypotheses must be
constructed after this extra type is killed; they are not supplied by
this theorem.
\end{theorem}

\begin{proof}
\textbf{1. The prime and the character.}
Trial division by the primes at most $\sqrt{3229}<57$ proves that
$3229$ is prime.  Moreover
\[
 q-1=12\cdot269,\qquad q\equiv1\pmod3,
 \qquad \operatorname{ord}_{\mathbf F_q^\times}(6)=3228.
 \tag{G119.7}
\]
The last equality follows by checking the three prime divisors
$2,3,269$ of $3228$.
Hence $\theta$ is defined and has exact order $269$ and conductor $q$.
It is even, and $\eta\theta$ is primitive odd of conductor $N=3q$.
The field $T$ is real because $\theta(-1)=1$.
All elementary computations, including the discrete logarithm of
$269$ modulo $3229$ modulo $269$, are replayed in the accompanying
integer certificate; no class-group computation is used.

\textbf{2. An integral Eisenstein form with cuspidal reduction.}
Choose a finite extension of $\mathbf Q_p$ containing the character
values and the prime-to-$p$ roots of unity needed for full level $4N$,
and write $\mathcal O$, $\varpi$, and $\kappa$ for its integer ring,
uniformizer, and residue field.
This temporary coefficient enlargement is part of the existence
construction; no previously fixed deformation coefficient ring is
replaced by it.
Set $\chi=\eta\theta$ and consider
\[
 E_k(1,\chi)
  =\frac{L(1-k,\chi)}2+
    \sum_{n\ge1}\left(\sum_{d\mid n}\chi(d)d^{k-1}\right)Q^n.
 \tag{G119.8}
\]
Here $Q$ is the Fourier variable, not the auxiliary prime.

For completeness the constant-term integrality does not divide by
$k$.  In Katz's Eisenstein measure use the test function
$x^{k-1}\chi(t)$ on $\mathbf Z_p\times\mathbf Z/N\mathbf Z$ and
the smoothing pair $(a,b)=(2,2)$.
Theorem 3.3.3 and Corollary 3.3.7 give the integral moment
\[
 2(1-2^k\chi(2))E_k(1,\chi).
 \tag{G119.9}
\]
The same formula applies to the imprimitive function
$\chi_0(t)=\eta(t)$ if $q\nmid t$ and $0$ otherwise.
The functions $\chi$ and $\chi_0$ are congruent modulo $\varpi$.
Since $\eta(2)=-1$ and $2^{134}=-1$ in $\mathbf F_{269}$,
$1+2^k$ is nonzero for every odd $k$: an odd power of $2$ cannot
equal $-1$, because $2$ has order $268$.
The factors in (G119.9) are therefore units and have the same
reduction for the two functions.
Continuity and $\mathcal O$-linearity of the integral measure give
\[
 \frac{L(1-k,\chi)}2\equiv
 \frac{L(1-k,\chi_0)}2
  =(1-\eta(q)q^{k-1})\frac{L(1-k,\eta)}2
  \equiv0\pmod{\varpi}.
 \tag{G119.10}
\]
The same smoothing argument shows integrality of the last
$L$-value; thus the displayed Euler factor is used between integral
quantities.  The argument applies also when $p\mid k$.

The q-expansion principle makes (G119.8) an integral holomorphic
section on the fine modular curve used below.
The complete constant formula of Dasgupta--Kakde, Theorem 4.5(1),
specialized to the primitive second character $\chi$ of conductor
$N$ and first character $1$, says that the constant at a cusp is
zero unless that cusp lies over the infinity class $C_\infty(N,N)$.
At each remaining cusp it is a root of unity times
$L(1-k,\chi)/2$.
Indeed both Gauss sums in their quotient are the same, the
conductor ratio is $1$, and the extra Euler-product set is empty.
For a rational cusp represented by a primitive integral pair
$(a,c)$, their normalization has no nonunit factor.
Pullback to full level $4$ changes only the Fourier parameter and
the trivialization by prime-to-$p$ factors.
Consequently (G119.10) kills the constant at every cusp.
The reduction $\overline E$ is a nonzero cusp section, since its
coefficient at $Q$ is $1$.

\textbf{3. The entire compact diamond cover is free in order $p$.}
Let $Y$ be the proper fine modular curve for
$\Gamma_1(N)\cap\Gamma(4)$, on a fixed full-level-$4$ pairing
component over $\mathcal O$.
It is smooth, with reduced cusp divisor $C_Y$ and Hodge line
$\omega_Y$.  One construction starts with the full level $4N$
smooth model and quotients by the upper-unipotent groups at $3$
and $q$, of order $N$, which is invertible in $\mathcal O$.
On the open moduli scheme these groups act freely; at a cusp the
quotient of its Tate parameter has the form $t\mapsto t^d$ with
$d\mid N$.  Thus the quotient is smooth and its cusps remain
finite etale over $\mathcal O$.
Katz's full-level cusp descriptions and the Kodaira--Spencer
isomorphism give
\[
 \omega_Y^{\otimes2}\simeq\Omega^1_{Y/\mathcal O}(C_Y).
 \tag{G119.11}
\]
The isomorphism is natural for the diamond action.

Let $D=(\mathbf Z/N\mathbf Z)^\times$ act by diamonds and let
$P\subset D$ be its $p$-Sylow subgroup.  By (G119.7), $|P|=p$.
On the open curve a fixed diamond point would give an automorphism
of the elliptic curve fixing the full level-$4$ structure.
Such an automorphism is the identity, so it cannot multiply the
order-$N$ point by a nonidentity element of $P$.

Here is the cusp calculation, including the special fiber.
Write a rational cusp as $\gamma\infty$, where the first column
of $\gamma\in\mathrm{SL}_2(\mathbf Z)$ is $(a,c)^t$.
A parabolic stabilizer has, up to its harmless sign, the form
\[
 \gamma\begin{pmatrix}1&n\\0&1\end{pmatrix}\gamma^{-1}
 =\begin{pmatrix}1-acn&a^2n\\-c^2n&1+acn\end{pmatrix}.
 \tag{G119.12}
\]
Membership in $\Gamma_0(N)$ requires $N\mid c^2n$.
As $N$ is squarefree, this implies $N\mid cn$ and therefore the
diagonal diamond entries are $1$ modulo $N$.
Allowing a sign gives at most an element of order $2$, never a
nonidentity element of $P$.
The same cusp set and stabilizer calculation apply in every
geometric fiber: the full prime-to-$p$ level cusp scheme is finite
etale, its Tate level-structure labels are constant after the chosen
unramified extension, and the preceding quotient by the
prime-to-$p$ unipotent group takes the same orbits in all fibers.
No inference from generic freeness alone is being made.
Thus $P$ acts freely on every geometric point of $Y$.

The free finite constant-group quotient
\[
 \pi:Y\longrightarrow X=Y/P
 \tag{G119.13}
\]
is a finite etale $P$-torsor, including at the cusps.
The quotient exists as a scheme by the free-action
quotient theorem; properness and smoothness descend along this
finite etale cover.
The Hodge line and reduced cusp divisor descend to $\omega_X,C_X$,
and (G119.11) descends to
$\omega_X^2=\Omega^1_{X/\mathcal O}(C_X)$.

\textbf{4. Lift in the character line, without a $1/p$ idempotent.}
Faithfully flat descent of the trivial line on $Y$ with its
$P$-linearization $\theta^{-1}$ defines a line bundle
$\mathcal L_\theta$ on $X$ for which
\[
 H^0(X,\omega_X^k\mathcal L_\theta(-C_X))
  =H^0(Y,\omega_Y^k(-C_Y))[\theta].
 \tag{G119.14}
\]
The bracket denotes the exact $\theta$-eigencondition under the
standard diamond convention.  This is descent, not an averaged
isotypic projector.  The formula also defines the sign of the
linearization if diamonds are expressed by inverse pullback.
Since $\theta^p=1$, $\mathcal L_\theta^p\simeq\mathcal O_X$;
in particular it has degree zero on each geometric fiber.
Since $\theta\equiv1\pmod\varpi$, its reduction is the trivial
line with its canonical descent datum.

On a geometric special fiber Serre duality gives
\[
 H^1(X_\kappa,\omega_X^k\mathcal L_\theta(-C_X))^\vee
 =H^0(X_\kappa,\omega_X^{2-k}\mathcal L_\theta^{-1})=0.
 \tag{G119.15}
\]
Indeed $\deg\omega_X>0$ (pull back to the full-level modular
curve), while $2-k<0$ and the character line has degree zero.
Coherent cohomology is finite over $\mathcal O$.
The exact reduction sequence and Nakayama's lemma show that its
$H^1$ vanishes, and give the surjection
\[
 H^0(X,\omega_X^k\mathcal L_\theta(-C_X))
 \twoheadrightarrow
 H^0(X_\kappa,\omega_X^k(-C_X)).
 \tag{G119.16}
\]
Here the right side contains the descent of $\overline E$.

Lift it by (G119.16) and pull back to $Y$.
Write $D=P\times D'$; $|D'|$ is prime to $p$.
Apply the $\eta$-diamond projector for $D'$ and the projector
forgetting the auxiliary full level $4$.
The latter averages the commuting group
$\mathrm{SL}_2(\mathbf Z/4\mathbf Z)$ of order $48$.
All these denominators are $p$-adic units; the two actions commute
with $P$ by their disjoint level factors.
They fix $\overline E$, so the lifted reduction remains nonzero.
We have obtained a cusp form in the exact $\eta\theta$ character
space at level $N$ with reduction $\overline E$.
No division by $269$ has occurred.

\textbf{5. The eigenlift is ordinary and primitive.}
The integral cusp module just constructed is finite and
torsion-free over $\mathcal O$, hence free.
It is stable under all integral Hecke operators, including
$U_3,U_q$ and $T_p$.
By (G119.8), its nonzero reduction is a common eigenvector with
good-prime eigenvalues $1+\eta(r)r^{k-1}$ and with
$U_3=U_q=1$; in particular $T_p=1$ in the residue field.
The full Deligne--Serre lifting lemma, Lemma 6.11, applies to this
commuting set, and gives a characteristic-zero common eigenform
after a finite coefficient extension with these residual
eigenvalues.  Its coefficient at $Q$ can be normalized to $1$:
if that coefficient were zero, the identities
$a_1(T_n f)=a_n(f)$ for all $n$ would make every Fourier
coefficient zero.
The classical Hecke algebra is defined over a number field.
Its finitely many algebra generators have algebraic eigenvalues,
so the normalized eigenform has a number field of coefficients
and a place compatible with the chosen $p$-adic embedding.

The primitive newform below this eigensystem has conductor
dividing $N$, and its nebentype is the primitive character
$\eta\theta$ of conductor $N$.
It must therefore have conductor exactly $N$.
There are no oldform multiplicities left at this level, and the
normalized form is $g$ in (G119.2).
The operator $T_p$ was included in the lift, so $a_p(g)$ is a
unit.  This conclusion does not require an ordinary projector on
the Eisenstein module or a residual-irreducibility hypothesis.

\textbf{6. Local types and the square-root twist.}
Use the standard attached-representation theorem for classical
cuspidal newforms, its purity, and local--global compatibility.
These are the retained Deligne--Langlands--Carayol and Saito
inputs, not new assertions proved by the eigenlift.
At each of $3,q$ the conductor of $g$ has exponent $1$ and its
determinant is ramified with exponent $1$.
A two-dimensional Weil--Deligne representation with these
properties is a sum of one unramified character and one
conductor-one character, with $N=0$.
Indeed an irreducible Weil representation has no inertia
invariants and conductor at least $2$; a nonzero monodromy
operator forces the two inertia characters to be equal, which
gives either unramified determinant or conductor at least $2$.
Thus $R_g$ has inertia $1\oplus\eta$ at $3$ and
$1\oplus\Theta^{-1}$ at $q$, exactly, and is unramified away
$3qp$.

At $p$ the level is prime to $p$, so $R_g$ is crystalline, with
Hodge filtration indices $0,k-1$ and crystalline Frobenius
slopes $0,k-1$ since $a_p(g)$ is a unit.
The slope-zero line has Hodge number $0$ by weak admissibility;
it is admissible and yields an unramified subrepresentation.
The quotient has inertia $c_p^{1-k}$.
This is also the usual ordinary-newform local theorem.
The comparison and admissible-filtered-module equivalence are
retained upstream inputs; only this two-slope deduction is used.

Equation (G119.1) and the actual Hecke congruences give
$\overline R_g^{\rm ss}=1\oplus\eta\overline c_p^{-2e}$
by Chebotarev and semisimplification.
The primitive twist in (G119.4) corresponds to the inverse
Galois finite twist, so (G119.5) follows, not its opposite.
Taking determinants and reducing, using $\overline\Theta=1$
and $\psi^{-1}=\psi$, proves (G119.6).
At $q$ the twist gives the two characters
$\Theta^{1/2},\Theta^{-1/2}$, both nontrivial conductor-one
characters.  Hence the new conductor exponent there is $2$.
At $2$, $R_g$ and $\eta\Theta^{1/2}$ are unramified and
$\psi$ has conductor exponent $2$, so the scalar rank-two
twist has conductor exponent $4$.
At $3$ the $\eta$ twist just interchanges $1$ and $\eta$.
This proves the exact conductor $2^4\cdot3\cdot q^2$ of $f_0$,
the stated inertias, and $N=0$ everywhere away from $p$.
Its twist at $p$ is unramified before $c_p^e$, so ordinarity
persists and the inertia characters of $V$ are as asserted.
Absolute irreducibility is the usual cuspidal attached-representation
property, unchanged by scalar twists.

\textbf{7. What the cyclic field does and does not give.}
Since $\Theta$ has order $p$, the inertia at $q$ becomes trivial
on $G_T$; $T$ is unramified at every other finite prime.
Restriction of the absolutely irreducible rank-two $V$ to this
normal subgroup of index $p$ remains irreducible.
For example, Clifford semisimplicity and Frobenius reciprocity
give
\[
 \dim\operatorname{End}_{G_T}(V)
   =\sum_{\nu\in\widehat{\operatorname{Gal}(T/\mathbf Q)}}
       \dim\operatorname{Hom}_{G_{\mathbf Q}}(V,V\nu)=1.
 \tag{G119.17}
\]
A nontrivial term would imply $V\simeq V\nu$, hence
$\nu^2=1$ by determinants, impossible for an order-$269$
character.  Thus only the trivial term occurs.
The elementary calculation $[269]=[6]^{175}$ in
$\mathbf F_q^\times/(\mathbf F_q^\times)^p$
gives the stated Frobenius in $T$; no equality of those two
representatives in $\mathbf F_q^\times$ is asserted.
It has order $269$, so $T_p/\mathbf Q_p$ is unramified of
degree $269$.

Finally $e=267$, $\psi=\chi_{-4}$ obey the original parity and
residual separation conditions: $\psi(-1)=(-1)^e$ and
$268\nmid2e$.
However every $q$-inertia element lies in the original residual
kernel, whereas (G119.5) is nontrivial on that inertia.
Thus this auxiliary form supplies no original IsPt assertion and
contradicts none of AC-G116 or AC-G117.
The exact compact-type, coefficient, host, and numerical
comparisons needed after passage to $T$ remain separate.
\end{proof}

\paragraph{Geometric dictionary and source scope.}
The finite etale diamond cover (G119.13) and its rank-one
character descent replace an illegal order-$269$ averaging
operation.  The cusp divisor imposes vanishing at every boundary
point, and Serre duality proves actual integral base change in
that line.  The later field $T$ is the finite Artin cover that
kills the new tame character; this geometric dictionary does not
construct a unitary automorphic packet or prove a Selmer
vanishing statement.
The primary applied sources are Katz 1973, Sections 1.4--1.7;
Katz 1977, Section 2.5 and Theorem 3.3.3/Corollary 3.3.7;
Dasgupta--Kakde, Theorem 4.5(1) and its complete printed proof;
Deligne--Serre, Sections 1.6--1.9 and Lemma 6.11 with its complete
printed proof; and the free-quotient theorem, Stacks tag 07S7.
Local--global compatibility and the crystalline comparison are
the retained classical modular inputs, with the exact geometric
Frobenius convention checked against the full native Newform
declaration.  No Lean execution, complete new proof of those
deep comparison theorems, or original-model support is claimed.
